\documentclass[11pt,a4paper]{article}
\usepackage[margin=23mm]{geometry}
\usepackage{amsmath,amssymb,xcolor,hyperref}
\definecolor{ink}{RGB}{20,63,69}
\hypersetup{colorlinks=true,urlcolor=ink,pdfauthor={Lee Crossley},pdftitle={Combined prime contributions in Johnston's sign problem}}
\setlength{\parindent}{0pt}
\setlength{\parskip}{5pt plus 1pt minus 1pt}
\setlength{\emergencystretch}{2em}
\newcommand{\head}[1]{\section*{\color{ink}#1}}
\newcommand{\li}{\operatorname{li}}
\raggedbottom
\begin{document}
\begin{center}
{\Large\bfseries Combined prime contributions in Johnston's sign problem\par}
\vspace{6pt}
{Lee Crossley \quad$\cdot$\quad 21 November 2026}
\end{center}
Fix a prime \(q\ge11\), put \(h=q-1\), and let \(a\) be a reduced class.
The normalized integrated prime-counting error is
\[
 E_{q,a}(x)=\frac{\log x}{x^{3/2}}
 \int_2^x\!\bigl(h\pi(t;q,a)-\li(t)\bigr)\,dt .
\]
We examine whether several primes can force \(E_{q,a}\) to take both signs
arbitrarily far out in every class. This note supplies an exact test, an
example where twelve terms are necessary, and a quantitative obstruction to
uniformly bounding the number of terms. The all-prime, all-class problem
remains open in this work. \href{https://arxiv.org/pdf/2407.03530#page=9}{Johnston's Problem 23.}

\head{The exact test for a collection of primes}
For a finite set of prime powers at positions \(n_i\) with positive weights \(w_i\),
write \(W=\sum_iw_i\) and \(S=\sum_iw_i n_i\). The line \(Wy-S\) is a global
minorant of \(\sum_iw_i(y-n_i)_+\), where \(t_+=\max(t,0)\). Completing the square gives
\[
 Wy-S-y^2/2=W^2/2-S-(y-W)^2/2 .
\]
Thus a collection with \(W^2>2S\) supplies a positive surplus at \(y=W\).
Conversely, if the full clipped sum exceeds \(y^2/2\) at any \(y\), its
active prefix of terms \(n_i\le y\) satisfies \(W^2>2S\). This is an
equivalence over prefixes, allowing any number of terms. The main-term test
alone omits the lower-order corrections in the actual prime-counting error.

Under Dirichlet GRH write \(E_{q,a}(x)=F_a(\log x)+o_q(1)\). The absolutely
convergent zero expansion makes \(F_a\) uniformly almost periodic. Every
strict profile value recurs, so two signed witnesses suffice. Use \(a\ne\pm1\),
\(b=a^{-1}\bmod q\), and \(B=\frac23\#\{r:r^2\equiv a\bmod q\}\). Set
\[
 \begin{aligned}
 w_a(n)&=h\,1_{n\equiv a(q)}\Lambda(n)+1_{n=q^j,\ j\ge1}\log q,\\
 D_b(s)&=\sum_nw_b(n)n^{-s},\qquad c_b=\mathrm{FP}_{s=1}D_b(s),\quad d_b=D_b(2).
 \end{aligned}
\]
The last series defines \(D_b\) first in \(\Re s>1\). Its continuation comes
from the character-weighted logarithmic derivatives of the primitive factors.
The forward Riesz identity from the preceding research note is
\[
 G(y):=y^{3/2}F_a(\log y)=\sum_n w_a(n)(y-n)_+-T(y),
\]
\[
 T(y)=y^2/2+B y^{3/2}+y\log y-y(1+\log q+c_b)+\log q+d_b,\qquad y\ge1 .
\]
Here \(T''(y)=1+3B/(4\sqrt y)+1/y>0\). On any finite search interval,
positive values exist exactly when some prefix satisfies
\(\max_y\{Wy-T(y)\}>S\), where the maximum is restricted to that interval.
The minorant argument works even when the maximizing \(y\) lies outside the
prefix's interval of activity. \href{https://github.com/leecrossley/integrated-prime-errors-mod-11/blob/v2.0.0/range/PROOF.md}{Fixed-modulus profile and recurrence.}

\newpage
\head{Twelve primes are necessary in one class}
At \(q=10\,009,\ a=1\,258\) (a nonsquare), the first twelve primes in the class are
\[
\begin{gathered}
251483,\ 271501,\ 311537,\ 451663,\ 551753,\ 791969,\\
872041,\ 972131,\ 1052203,\ 1212347,\ 1292419,\ 1412527 .
\end{gathered}
\]
At \(y=1\,611\,211\), these are all the progression terms and
\(F_a(\log y)>1.00540335\). For this forward identity twelve is the exact
minimum number of selected progression terms for a positive test at any
real \(y\ge1\), even when every common \(q\)-power correction is retained.

Indeed \(b=1\,456\) and interval calculation gives
\(1+\log q+c_b<-2.1995\). The common terms add at most \(y\log y\) at every
\(y\ge1\). Since \(B\ge0\) and \(d_b>0\), a positive corrected test requires
the chosen progression terms alone to exceed \(y^2/2\). For at most twelve
terms this requires \(y<2h\cdot12\log y\), hence \(y<6\,000\,000\), using
monotonicity beyond this endpoint. Direct sieving finds 40 ordinary primes
and no higher powers in the class up to that cutoff.

For fixed \(y\), the function \(t\mapsto (y-t)\log t\), \(2\le t\le y\), is
strictly concave. Its derivative is \(y/t-1-\log t\). The largest \(k\)
active weights in the ordered list of primes therefore occur in a consecutive
block, with ties resolved at its two ends. All 385 consecutive blocks of
length one through eleven have \(W^2/2<S\). Completing the square excludes
them at every scale. A 160-bit interval calculation verifies the positive
corrected value for the displayed list. The integer sieve for this example
is independent of the progression sieve in the continuous search.

\head{A finite search over every real scale}
For the reflected profile put \(H(y)=F_a(-\log y)/\sqrt y\). If the active
inverse-class prefix has \(P=\sum_{n\le y}w_b(n)/n\) and
\(Q=\sum_{n\le y}w_b(n)/n^2\), the reflected Riesz identity gives
\[
 H(y)=1+\log y+c_b-P+y(Q-d_b)-1/(2y)-B/\sqrt y .
\]
Both \(G\) and \(H\) are strictly concave between successive prime-power
positions in their respective progressions, including the common \(q\)-powers.
In particular \(H''(y)=-y^{-2}-y^{-3}-3B/(4y^{5/2})\) on each piece.
At any interior rational \(z\), or at an endpoint using the piece's one-sided
derivative, the tangent line majorizes that whole piece. If the tangent is
strictly negative at both endpoints then the profile is negative everywhere
on that piece. Conversely, a strictly positive interval value at \(z\)
certifies a positive witness. Floating point proposes \(z\). Each conclusion
uses outward-rounded ball arithmetic on the exact dyadic rational proposed.

\newpage
\head{What happens as the search interval grows}
At \(q=10\,009\) an earlier geometric grid found positive profiles for
9,779 of the 10,006 classes with \(a\ne\pm1\), and negative profiles for
all 10,006. The continuous search starts with the other 227 and retains
every progression term. Each row exhausts both directions on \(1\le y\le Y\)
for every class left in its last column.
\begin{center}
\begin{tabular}{r r r r}
\hline
\(Y\)&Classes entering&New positive witnesses&Negative throughout\\
\hline
\(2\cdot10^6\)&227&6&221\\
\(2\cdot10^8\)&221&198&23\\
\(10^{10}\)&23&18&5\\
\(10^{11}\)&5&3&2\\
\(10^{12}\)&2&0&2\\
\hline
\end{tabular}
\end{center}
At \(a=148\) and \(a=9\,100\) this proves \(F_a(u)<0\) for the full real
interval \(|u|\le\log(10^{12})\). This assertion concerns the limiting
profile. Both classes are quadratic nonresidues. All 5,002 quadratic residues
other than \(\pm1\) already have a positive arithmetic witness by \(y=10^{11}\).
The separately verified low-zero certificate at this modulus also gives
\(\limsup E_{q,a}>1/4\) for both classes. Hence their positive profile values occur
outside the tested interval. A finite profile search cannot substitute for
a theorem about arbitrarily large prime moduli.

At 128 and 160 bits the programs visited the same 18,578,682 pieces across all
five stages, with the same dyadic positive witnesses and the same surviving
classes. Some scale intervals recur across stages. The previous low-zero
certificate used two interval methods, which share FLINT's endpoint sign tests. The continuous search
uses the same algorithm at both precisions. \href{https://arxiv.org/pdf/1305.3087#page=14}{Platt's finite nonvanishing input.}

\head{How many terms can a class require?}
At a position \(n\) with weight \(h\Lambda(n)\), one has
\[
 \sup_{y>0}\frac{h\Lambda(n)(y-n)_+}{y^2}
 =\frac{h\Lambda(n)}{4n}.
\]
If a progression has no prime power up to \(Y>e\), any \(k\) progression
terms contribute at most \(kh\log Y/(4Y)\) after division by \(y^2\).
Even the best choice of terms and scale therefore needs
\(k>2Y/(h\log Y)\) to exceed the main term \(y^2/2\).
For ordinary primes alone the theorem of Li, Pratt and Shakan already
implies that no fixed \(k\) works this way for every class at every large
prime. The next page handles prime powers and the inverse class simultaneously.
\href{https://arxiv.org/pdf/1607.02543#page=2}{Li, Pratt and Shakan, Theorem 1.1.}

\newpage
\head{Long gaps for a residue and its inverse}
Let \(\log_j\) denote the \(j\)-fold iterated logarithm, and put
\(L(q)=(\log_2q)(\log_4q)/\log_3q\). For every fixed \(0<\varepsilon<1\),
all sufficiently large primes \(q\) admit at least \(q^{1-\varepsilon}\)
residues \(a\ne\pm1\) such that neither \(a\) nor \(a^{-1}\) contains a
prime power up to
\[
 Y=c_\varepsilon q\log q\,L(q)
\]
for a positive constant \(c_\varepsilon\). This arithmetic assertion is
unconditional. Thus in both progressions a main-term-dominating selection
requires \(\gg_\varepsilon L(q)\) terms even with access to arbitrarily
large primes and prime powers.

Take \(\tau=\min(\varepsilon/8,1/16)\), \(u=\tau\log q\) and let
\(m=\prod_{p\le u}p=q^{\tau+o(1)}\). The covering estimate of Ford, Green,
Konyagin, Maynard and Tao supplies an integer \(v\) and a length
\(J\asymp_\varepsilon\log q\,L(q)\) such that every \(v+j,\ 0\le j<J\),
has a prime factor dividing \(m\). We truncate a longer cover if necessary.
The set \(S=\{a:1\le a<q,\ a\equiv qv\bmod m\}\) has size \(q/m+O(1)\).
For \(a\in S\), each \(a+jq,\ 0\le j<J\), has a prime factor \(\le u\).
Any prime power among them must therefore be a power of a prime \(\le u\),
including exponent one. Across all \(a\in S\) there are at most
\(O(\pi(u)\log(q(J+1)))=O(\log^2q)\) bad residues. Every member of an
eligible progression up to \(qJ\) is among the tested \(J\) positions.
\href{https://arxiv.org/pdf/1412.5029#page=3}{Ford et al., equation (1.2).}

Write \(a=r+ms,\ b=r+mt\) for the same least representative \(r\bmod m\).
After division by \(m^2\bmod q\), \(ab\equiv1\bmod q\) becomes
\((s+r/m)(t+r/m)\equiv m^{-2}\bmod q\). The two variables range over
translated intervals of the same length \(|S|\) modulo \(q\).
The modular-hyperbola count is uniformly
\[
 \#\{a\in S:a^{-1}\in S\}=\frac{|S|^2}{q}
              +O(\sqrt q\log^2q)=q^{1-2\tau+o(1)} .
\]
In fact each quadratic type contains at least \(q^{1-\varepsilon}\) such residues.
Let \(\chi\) be the Legendre symbol and \(e_q(t)=\exp(2\pi i t/q)\).
The normalized Fourier coefficients of \(1_S\) have total absolute value
\(O(\log q)\) by the geometric-series bound for an interval and dilation by \(m\).
Expanding both indicators in \(\sum_a\chi(a)1_S(a)1_S(a^{-1})\)
gives a discrepancy \(O(\sqrt q\log^2q)\). Indeed the complete sums
\(\sum_{a\ne0}\chi(a)e_q(ua+v/a)\) have modulus at most \(2\sqrt q\).
For \(v\ne0\) the Sali\'e formula gives
\(\tau(\chi)\chi(v)\sum_{z^2=4uv} e_q(z)\), and for \(v=0\) the
Gauss bound suffices. Half the unweighted count therefore survives for
each type up to that discrepancy. Removing the bad classes, their inverses
and \(\pm1\) costs \(O(\log^2q)\). Since \(2\tau<\varepsilon\), the
claim follows after reducing \(c_\varepsilon\).
For example, with \(q=1\,000\,003,\ m=210,\ v=2,\ J=9\), an independent
sieve finds four squares and six nonsquares, including the pair \(1646,318956\).
\href{https://doi.org/10.4064/aa142-1-5}{Chan and Shparlinski, equation (1).}
\href{https://www.math.ucdavis.edu/~junxian/paper/DivisorTichmarsh5.pdf\#page=18}{Li, Lemma 6.3.}

Under Dirichlet GRH the preceding note gives
\(c_b=h\Lambda(b)/b+O(\log^2q)\). Here \(\Lambda(b)=0\). Up to \(Y\) the
only terms in the forward identity are the common \(q\)-powers. Hence
\(\sqrt y\,F_a(\log y)\le-y/2+1+\log q+c_b\le-y/3\)
for \(C\log^2q\le y\le Y\) with an absolute \(C\), and likewise for the
inverse. This blocks every early forward prime-sum test on that whole
interval. Later contributions, and the two-sided spectral method, remain
available. The covering bound and inverse-count estimate are established
results. The paired prime-power application and the numerical examples are
recorded here without a claim of historical priority.
\href{https://arxiv.org/pdf/0912.4908#page=31}{Fiorilli and Martin, Propositions 3.10--3.11.}
\end{document}
